\chapter{Допустимые деформации и устойчивость режимов}

Структурная классификация имеет смысл только тогда, когда можно отличить
изменение одной реализации от перехода к иной сущности. Для этого вводится
класс допустимых деформаций $\Delta_{\mathrm{adm}}$.

\section{Слабые деформации}

Деформация $\delta$ называется слабой, если она сохраняет инвариантное ядро и
тип редукции:
\begin{equation}
 \mathcal I(\delta\mathfrak S)\simeq\mathcal I(\mathfrak S),
 \qquad
 \mathcal R_{\mathrm{type}}(\delta\mathfrak S)=
 \mathcal R_{\mathrm{type}}(\mathfrak S).
\end{equation}

К слабым изменениям могут относиться смена координат, базиса, допустимой
параметризации или масштаба чтения. Они не должны создавать новые физические
степени свободы только за счёт переописания.

\section{Сильный разрыв}

Если деформация меняет инвариантное ядро, число независимых секторов или класс
глобальной допустимости, она рассматривается как переход к иной структуре.
Такой переход нельзя использовать как обычную проверку устойчивости прежней
модели.

\section{Устойчивый режим}

Режим $r$ считается устойчивым в области $U$, если для всех достаточно малых
$\delta\in\Delta_{\mathrm{adm}}$ выполняются
\begin{equation}
 r(\delta O)=r(O),
 \qquad
 d_{\mathfrak R}\bigl(\mathfrak R(\delta O),\mathfrak R(O)\bigr)<\varepsilon.
\end{equation}

Первое условие сохраняет тип, второе ограничивает движение внутри области.
Граница устойчивости определяется заранее выбранным $\varepsilon$, а не после
получения удобного результата.

\section{Запрещённые подмены}

Нельзя выдавать за устойчивость:
\begin{itemize}
 \item сохранение результата после перенастройки свободного коэффициента;
 \item выбор новой ветви после сравнения с целевыми данными;
 \item замену физического изменения сменой обозначений;
 \item исключение неудобного сектора без нового критерия допустимости.
\end{itemize}

\chapter{Структурная динамика}

Карта режимов становится динамической, если допустимой деформации сопоставлена
траектория
\begin{equation}
 t\longmapsto\mathfrak R_t(O)=
 \bigl(\Xi_t(O),\Upsilon_t(O)\bigr).
\end{equation}
Параметр $t$ пока обозначает порядок деформации или масштаб чтения, а не
обязательно физическое время.

\section{Поле структурного потока}

Локальное движение на карте задаётся полем
\begin{equation}
 \mathbf V(\Xi,\Upsilon)=
 \left(\frac{d\Xi}{dt},\frac{d\Upsilon}{dt}\right).
\end{equation}

Первая компонента показывает усиление или ослабление некаркасного вклада,
вторая --- смещение между фазовой и внутренней ориентацией.

\section{Неподвижные режимы}

Точка $\mathfrak R_*$ является неподвижной, если
\begin{equation}
 \mathbf V(\mathfrak R_*)=0.
\end{equation}
Её физическое толкование допустимо только после вывода потока из конкретного
действия или закона обновления. Нуль произвольно выбранного вспомогательного
поля не является доказательством устойчивой частицы или вакуума.

\section{Переходы между областями}

Траектория может пересекать границу каркасного, фазового, внутреннего или
смешанного режима. Для каждого пересечения требуется определить:
\begin{enumerate}
 \item сохраняется ли инвариантное ядро;
 \item непрерывен ли операторный носитель;
 \item меняется ли ранг или топологический класс;
 \item существует ли конечная стоимость перехода;
 \item является ли конечный режим устойчивым.
\end{enumerate}

Без этих проверок карта описывает лишь изменение классификационных меток.

\section{Граница с физической динамикой}

Структурный поток становится физическим только после появления единого
временного закона, энергетической нормировки и наблюдаемой скорости. До этого
он служит языком организации деформаций и не предсказывает длительность или
вероятность процесса.

\chapter{Редукционное сжатие аппарата}

Развитая программа неизбежно вводит больше понятий, чем требуется её
минимальному основанию. Поэтому каждый этап расширения должен завершаться
сжатием.

\section{Расширенный аппарат}

Текущий рабочий набор можно представить как
\begin{equation}
 \mathfrak A_{\mathrm{ext}}=
 \bigl(\mathfrak S,\mathcal I,\mathcal R,
 \Xi,\Upsilon,\mathfrak O_{\mathrm{cls}},
 \mathcal Q_{\mathrm{sel}}\bigr).
\end{equation}

Здесь $\mathfrak O_{\mathrm{cls}}$ обозначает классы наблюдаемых, а остальные
компоненты фиксируют ядро, режимы и выбор моделей.

\section{Четыре шага сжатия}

Для каждого элемента выполняются:
\begin{enumerate}
 \item указание его несущей функции;
 \item пробное исключение;
 \item попытка восстановить функцию более короткими средствами;
 \item присвоение статуса.
\end{enumerate}

Допустимы три статуса: фундаментальный, выводимый и неопределённый. Полезность
в вычислении сама по себе не делает элемент фундаментальным.

\section{Что обязано сохраниться}

После сжатия должны сохраняться:
\begin{itemize}
 \item идентичность $\mathfrak S$;
 \item различимость режимов;
 \item реконструкция наблюдаемых;
 \item нетривиальный выбор между моделями;
 \item уже полученные условия провала.
\end{itemize}

Если более короткая схема теряет хотя бы одну из этих функций, сжатие было
чрезмерным.

\section{Научный смысл сжатия}

Цель процедуры состоит не в сокращении числа страниц, а в отделении предмета
теории от временных средств его поиска. Зрелая программа должна объяснять
больше меньшим числом независимых предположений и сохранять различающую силу.

Следующий крупный блок должен объединить пробный цикл применения протокола,
иерархию моделей и критерии перехода от методической программы к строгой
теории.